\section{Praktikum}
\label{sec:praktikum-bab11}

\begin{praktikumbox}
\vspace{1.1em}
\textbf{Judul:} Manipulasi bit dengan SHR, TEST, dan XOR.\\[0.3em]
\textbf{Sub-CPMK:} 3.2 \quad\textbullet\quad
\textbf{Alokasi waktu:} 100 menit \quad\textbullet\quad
\textbf{Bentuk:} individu
\end{praktikumbox}

\subsection{Tujuan Praktikum}

Setelah menyelesaikan praktikum ini, mahasiswa mampu menggunakan
\asm{SHR}, \asm{TEST}, dan \asm{XOR} untuk memeriksa serta mengubah
bit, lalu menjelaskan hasilnya berdasarkan isi register dan
\eng{flag}.

\subsection{Alat dan Bahan}

\begin{itemize}
    \item Komputer dengan emu8086 terpasang.
    \item DOSBox dengan TASM, TLINK, dan DEBUG.
    \item Berkas \file{paritas.asm} dan \file{xor-cipher.asm} dari
          \file{code/asm/}.
    \item Lembar pengamatan register dan \eng{flag}.
\end{itemize}

\subsection{Langkah Kerja}

\begin{enumerate}
    \item \textbf{Menguji program paritas.}
    \begin{enumerate}
        \item Rakit \file{paritas.asm} dengan TASM, lalu link
              dengan TLINK.
        \item Jalankan program dan catat keluarannya.
        \item Buka program pada emu8086, lalu jalankan
              \eng{single step} sampai perulangan selesai.
        \item Catat \reg{AL}, \reg{BL}, \reg{CX}, dan \flag{CF}
              pada setiap pergeseran.
        \item Jelaskan mengapa \reg{BL} berisi jumlah bit 1.
    \end{enumerate}

    \item \textbf{Mengubah paritas.}
    \begin{enumerate}
        \item Ubah \code{nilai} menjadi \bin{10110111} yang memiliki
              enam bit 1.
        \item Rakit ulang dan jalankan program.
        \item Pastikan keluaran berubah menjadi \code{GENAP}.
    \end{enumerate}

    \item \textbf{Menguji XOR cipher.}
    \begin{enumerate}
        \item Rakit dan jalankan \file{xor-cipher.asm}.
        \item Catat teks asli, hasil enkripsi, dan hasil dekripsi.
        \item Jelaskan mengapa hasil dekripsi sama dengan teks asli.
    \end{enumerate}

    \item \textbf{Mengubah kunci XOR.}
    \begin{enumerate}
        \item Ubah \code{kunci} dari \hex{20} menjadi \hex{01}.
        \item Rakit ulang, lalu jalankan program.
        \item Catat hasil enkripsi dan dekripsi.
        \item Jelaskan mengapa hasil enkripsi tidak lagi terbaca
              sebagai huruf.
    \end{enumerate}

    \item \textbf{Menguji mask.}
    \begin{enumerate}
        \item Tulis potongan program yang menyalakan bit 3 pada
              \reg{AL} = \hex{A5}, lalu memadamkannya kembali.
        \item Gunakan \asm{OR} dan \asm{AND} dengan mask yang tepat.
        \item Periksa hasilnya dengan \eng{single step}.
    \end{enumerate}

    \item \textbf{Memeriksa dengan DEBUG.}
    \begin{enumerate}
        \item Jalankan \file{debug paritas.exe}.
        \item Ketik \file{r} untuk melihat register.
        \item Ketik \file{t} beberapa kali dan catat perubahan
              \reg{AL}, \reg{BL}, dan \flag{CF}.
    \end{enumerate}
\end{enumerate}

\subsection{Hasil yang Dikumpulkan}

\begin{enumerate}
    \item Berkas \file{paritas.asm} dan \file{xor-cipher.asm} versi
          modifikasi.
    \item Tangkapan layar keluaran kedua program.
    \item Tabel pengamatan \reg{AL}, \reg{BL}, \reg{CX}, dan
          \flag{CF} pada perulangan paritas.
    \item Penjelasan sifat $A \oplus K \oplus K = A$.
    \item Catatan hasil percobaan mask.
\end{enumerate}

\begin{asesmenbox}
\vspace{1.1em}
\textbf{Asesmen praktikum 9 (bobot 5\% dari komponen praktikum).}

\begin{tabular}{@{}p{6.2cm}cp{4.6cm}@{}}
    \toprule
    \textbf{Indikator} & \textbf{Bobot} & \textbf{Kriteria} \\
    \midrule
    Ketepatan pengamatan pergeseran & 30\% &
    \reg{AL}, \reg{BL}, \reg{CF} tercatat benar \\
    \addlinespace
    Ketepatan analisis paritas & 25\% &
    Hubungan jumlah bit 1 dan \asm{TEST} dijelaskan \\
    \addlinespace
    Ketepatan analisis XOR & 25\% &
    Enkripsi dan dekripsi dijelaskan dengan sifat XOR \\
    \addlinespace
    Ketepatan penggunaan mask & 20\% &
    Mask dan instruksi yang dipilih benar \\
    \bottomrule
\end{tabular}
\end{asesmenbox}
